\documentclass[10pt,aspectratio=169]{beamer}
\input{../../_en_preambulo_beamer}
% Comandos y macros estadísticas compartidas para las presentaciones Beamer en inglés (Metropolis)

% Conjuntos numéricos básicos
\newcommand{\R}{\mathbb{R}}
\newcommand{\N}{\mathbb{N}}
\newcommand{\Q}{\mathbb{Q}}
\newcommand{\Z}{\mathbb{Z}}
\newcommand{\set}[1]{\left\{ #1 \right\}}
\newcommand{\abs}[1]{\left|#1\right|}
\newcommand{\norm}[1]{\left\| #1 \right\|}

% Comandos estadísticos y probabilísticos del curso
\newcommand{\Var}{\operatorname{Var}}
\newcommand{\cov}{\operatorname{Cov}}
\newcommand{\corr}{\rho}
\newcommand{\comb}[2]{\begin{pmatrix} #1 \\ #2 \end{pmatrix}}
\newcommand{\s}{\sigma}
\newcommand{\particion}{\mathfrak{P}}
\newcommand{\card}[1]{\#\left( #1 \right)}
\newcommand{\seq}[3]{\set{#1}_{#2}^{#3}}
\newcommand{\E}{\mathbb{E}}
\newcommand{\Prob}{\mathbb{P}}

% Entornos pedagógicos y cajas en inglés académico natural (soportando tanto nombres en inglés como en español)
\newenvironment{discussion}{%
  \begin{alertblock}{Discussion Question}%
}{%
  \end{alertblock}%
}
\newenvironment{discusion}{%
  \begin{alertblock}{Discussion Question}%
}{%
  \end{alertblock}%
}

\newenvironment{reflection}{%
  \begin{alertblock}{Classroom Reflection}%
}{%
  \end{alertblock}%
}
\newenvironment{reflexion}{%
  \begin{alertblock}{Classroom Reflection}%
}{%
  \end{alertblock}%
}

\newenvironment{paradox}{%
  \begin{alertblock}{Exploratory Paradox}%
}{%
  \end{alertblock}%
}
\newenvironment{paradoja}{%
  \begin{alertblock}{Exploratory Paradox}%
}{%
  \end{alertblock}%
}

\newenvironment{motivation}{%
  \begin{exampleblock}{Motivation: Data Science in Practice}%
}{%
  \end{exampleblock}%
}
\newenvironment{motivacion}{%
  \begin{exampleblock}{Motivation: Data Science in Practice}%
}{%
  \end{exampleblock}%
}

\newenvironment{quickchallenge}{%
  \begin{block}{Quick Team Challenge}%
}{%
  \end{block}%
}
\newenvironment{retoclase}{%
  \begin{block}{Quick Team Challenge}%
}{%
  \end{block}%
}


\title[Foundations and Axioms]{Section 2.3: Probability Foundations and Axioms}
\subtitle{Formalizing Randomness: From Frequentist Counts to Kolmogorov Axioms}
\author[J. Castillo Colmenares]{Juliho Castillo Colmenares}
\institute{Tecnológico de Monterrey}
\date{}

\begin{document}

% Slide 1: Title Page
\begin{frame}[plain]
  \titlepage
\end{frame}

% Slide 2: Roadmap and Position Indicator
\begin{frame}{Roadmap — Chapter 02: Probability Theory}
  Where are we in our exploration of probability theory? \pause
  \begin{itemize}
    \item \textbf{02.01 Introduction to probability} \emph{(Completed)} \pause
    \item \textbf{02.02 Sets and partitions notation} \emph{(Completed)} \pause
    \item \color{TecRojo}\textbf{02.03 Probability foundations and axioms \emph{(Today)}}\color{black} \pause
    \item \textbf{02.04 Conditional probability and independence} \emph{(Next session)} \pause
    \item \textbf{02.05 Bayes' theorem} \pause
    \item \textbf{02.06 Random sampling}
  \end{itemize}
  \vspace{0.6em}
  \begin{block}{Section 2.3 Objective}
    Establish the three Kolmogorov axioms as the rigorous mathematical foundation for quantifying uncertainty in discrete and continuous experiments, deriving the fundamental rules of probability calculus.
  \end{block}
\end{frame}

% Slide 3: Motivation
\begin{frame}{The Logical Consistency Problem in Data Science}
  \begin{motivacion}[Why do we need axioms for uncertainty?]
    If a medical diagnostic algorithm assigns a $60\%$ probability that a patient has Type 1 diabetes and a $55\%$ probability of Type 2 diabetes (being mutually exclusive), their sum yields $115\%$. An algorithmic model without strict axiomatic rules produces mathematical contradictions and absurd inferences.
  \end{motivacion}
  
  \vspace{0.4em}
  \pause
  \begin{itemize}
    \item \textbf{Quantifiable Uncertainty:} Probability is a measure function $P: \mathcal{C} \to \mathbb{R}$ that assigns a numerical value to the plausibility of an event. \pause
    \item \textbf{Universal Consistency:} Axioms guarantee that regardless of the experiment's nature (coins, genes, financial markets), the calculus remains logically infallible.
  \end{itemize}
\end{frame}

% Slide 4: Random Experiments and Sample Space
\begin{frame}{Random Experiments and Sample Space ($S$)}
  A \emph{random experiment} is a process whose specific outcome cannot be predicted with certainty in advance, yet whose set of all possible outcomes is completely known. \pause

  \vspace{0.3em}
  \begin{block}{The Sample Space ($S$)}
    The set $S$ consisting of all possible outcomes of a random experiment is called the \emph{sample space}, and each individual outcome is called a \emph{sample point}.
  \end{block}

  \vspace{0.3em}
  \pause
  \begin{itemize}
    \item \textbf{Tossing a coin:} $S = \{T, H\}$ or $\{0, 1\}$. \pause
    \item \textbf{Tossing two coins:} $S = \{HH, HT, TH, TT\}$. \pause
    \item \textbf{Quality Control (Bolt manufacturing):} $S = \{\text{defective}, \text{non-defective}\}$. \pause
    \item \textbf{Continuous Lifetime of a Component:} $S = \{t \in \mathbb{R} \mid 0 \leq t \leq T\}$, where $t$ is measured in operating hours.
  \end{itemize}
\end{frame}

% Slide 5: Types of Sample Space
\begin{frame}{Types of Sample Space: Discrete vs. Continuous}
  The mathematical structure of $S$ dictates which probabilistic tools must be used: \pause

  \vspace{0.4em}
  \begin{itemize}
    \item \textbf{Finite:} Contains a specific, countable number of sample points. Example: rolling a die, $S = \{1, 2, 3, 4, 5, 6\}$. \pause
    \item \textbf{Countable Infinite:} Contains as many points as the natural numbers $\mathbb{N}$ (can be listed one by one). Example: number of attempts until a server successfully responds, $S = \{1, 2, 3, \dots\}$. \pause
    \item \textbf{Uncountable Infinite:} Contains as many points as an interval on the real line $\mathbb{R}$. Example: API response latency in milliseconds, $S = (0, \infty)$.
  \end{itemize}

  \vspace{0.4em}
  \pause
  \begin{alertblock}{Operational Classification}
    If $S$ is finite or countably infinite, we call it a \textbf{discrete} space. If it is uncountably infinite, we call it a \textbf{continuous} space.
  \end{alertblock}
\end{frame}

% Slide 6: Events and Subset Algebra
\begin{frame}{Events and Subset Algebra}
  \small An \emph{event} $A$ is any subset of the sample space ($A \subset S$). If the observed outcome is an element of $A$, we say that \emph{event $A$ has occurred}. \pause

  \vspace{0.1em}
  \begin{itemize}\itemsep=0.1em
    \item \textbf{Simple or Elementary Event:} Consists of exactly one sample point (e.g., $\{HT\}$). \pause
    \item \textbf{Certain Event:} The entire sample space $S$ ($P(S) = 1$). \pause
    \item \textbf{Impossible Event:} The empty set $\emptyset$ ($P(\emptyset) = 0$).
  \end{itemize}

  \vspace{0.1em}
  \pause
  \begin{block}{Set Operations on Events}
    \begin{itemize}\itemsep=0.1em
      \item \textbf{Union ($A \cup B$):} Either $A$ occurs, or $B$ occurs, or both.
      \item \textbf{Intersection ($A \cap B$):} Both $A$ and $B$ occur simultaneously.
      \item \textbf{Complement ($A'$):} $A$ does not occur ($S \setminus A$).
      \item \textbf{Mutually Exclusive (Disjoint):} $A \cap B = \emptyset$.
    \end{itemize}
  \end{block}
\end{frame}

% Slide 7: Historical Approaches to Probability
\begin{frame}{The Three Classical Approaches to Probability}
  \footnotesize How do we assign the numerical probability value $P(A)$ to an event? \pause

  \vspace{0.1em}
  \begin{itemize}\itemsep=0.2em
    \item \textbf{1. Classical Approach (Laplace, 1812):} For $n$ \emph{equally likely} outcomes with $h$ favorable cases:
    $$P(A) = \frac{h}{n}$$ \pause
    \item \textbf{2. Frequentist / Empirical Approach (von Mises):} Repeating $n$ times under identical conditions:
    $$P(A) = \lim_{n \to \infty} \frac{h}{n}$$ \pause
    \item \textbf{3. Axiomatic Approach (Kolmogorov, 1933):} Defines probability mathematically via indisputable axioms, resolving equiprobability loops.
  \end{itemize}
\end{frame}

% Slide 8: Kolmogorov Axioms
\begin{frame}{The Three Kolmogorov Axioms (1933)}
  \small Let $S$ be a sample space and $\mathcal{C}$ the collection of events. A function $P: \mathcal{C} \to \mathbb{R}$ is a \emph{probability function} if it satisfies: \pause

  \vspace{0.2em}
  \begin{block}{Axiom 1 (Non-Negativity)}
    For any event $A \in \mathcal{C}$: $P(A) \geq 0$
  \end{block}

  \vspace{0.1em}
  \pause
  \begin{block}{Axiom 2 (Normalization or Certain Event)}
    The probability of the entire sample space $S$ is unity: $P(S) = 1$
  \end{block}

  \vspace{0.1em}
  \pause
  \begin{block}{Axiom 3 (Countable Additivity)}
    For mutually exclusive events $A_1, A_2, \dots$ ($A_i \cap A_j = \emptyset$ for $i \neq j$):
    $$P(A_1 \sqcup A_2 \sqcup \dots) = P(A_1) + P(A_2) + \dots$$
  \end{block}
\end{frame}

% Slide 9A: Fundamental Theorems Derivation (Part 1)
\begin{frame}{Fundamental Theorems Derived (1/2)}
  Using only the three axioms, we can rigorously prove all standard rules: \pause

  \vspace{0.4em}
  \begin{itemize}
    \item \textbf{Theorem 2.3.3 (Impossible Event):} $P(\emptyset) = 0$. \pause
    \vspace{0.3em}
    \item \textbf{Theorem 2.3.4 (Complement Rule):} For any event $A$:
    $$P(A') = 1 - P(A)$$ \pause
    \vspace{0.3em}
    \item \textbf{Theorem 2.3.7 (Binary Law of Total Probability):} For any binary partition of $B$ relative to $A$:
    $$P(A) = P(A \cap B) + P(A \cap B')$$
  \end{itemize}
\end{frame}

% Slide 9B: Fundamental Theorems Derivation (Part 2)
\begin{frame}{Fundamental Theorems Derived (2/2)}
  \begin{block}{Theorems 2.3.1 \& 2.3.2 (Monotonically Bounded Differences)}
    If $A \subset B$, then the probability of their set difference is exact:
    $$P(B \setminus A) = P(B) - P(A)$$
    By monotonicity, since $P(B \setminus A) \geq 0$, it follows directly that:
    $$P(A) \leq P(B)$$
    In particular, because $\emptyset \subset A \subset S$, every probability is bounded between $0\%$ and $100\%$:
    $$0 \leq P(A) \leq 1$$
  \end{block}
\end{frame}

% Slide 10: General Addition Rule (Inclusion-Exclusion)
\begin{frame}{Theorem 2.3.6: General Addition Rule}
  \small When two events $A$ and $B$ are \emph{not} mutually exclusive ($A \cap B \neq \emptyset$), simple addition counts their intersection twice: \pause

  \vspace{0.2em}
  \begin{block}{Addition Rule for Two Events}
    $$P(A \cup B) = P(A) + P(B) - P(A \cap B)$$
  \end{block}

  \vspace{0.2em}
  \pause
  \begin{block}{Addition Rule for Three Events (Inclusion-Exclusion)}
    \begin{align*}
      P(A \cup B \cup C) &= P(A) + P(B) + P(C) - P(A \cap B) \\
                         &\quad - P(B \cap C) - P(C \cap A) + P(A \cap B \cap C)
    \end{align*}
  \end{block}
\end{frame}

% Slide 11A: Python Laboratory (Part 1: Simulation)
\begin{frame}[fragile]{Python Lab: Frequentist vs. Classical}
  The script \texttt{02.03\_probability\_axioms.py} compares empirical vs. classical ($N = 10,000$):
  \vspace{0.2em}
  {\lstset{basicstyle=\fontsize{6.5pt}{7.5pt}\selectfont\ttfamily}
  \lstinputlisting[firstline=1, lastline=17]{../../code/02_teoria_probabilidad/02.03_probability_axioms.py}}
\end{frame}

% Slide 11B: Python Laboratory (Part 2: Axioms)
\begin{frame}[fragile]{Python Lab: Axiom Verification}
  Normalization and non-negativity validation on a 52-card deck:
  \vspace{0.2em}
  {\lstset{basicstyle=\fontsize{6.5pt}{7.5pt}\selectfont\ttfamily}
  \lstinputlisting[firstline=22, lastline=35]{../../code/02_teoria_probabilidad/02.03_probability_axioms.py}}
\end{frame}

% Slide 11C: Python Laboratory (Part 3: Addition Rule)
\begin{frame}[fragile]{Python Lab: Addition Rule}
  Inclusion-exclusion verification against direct counting:
  \vspace{0.2em}
  {\lstset{basicstyle=\fontsize{6.5pt}{7.5pt}\selectfont\ttfamily}
  \lstinputlisting[firstline=36, lastline=48]{../../code/02_teoria_probabilidad/02.03_probability_axioms.py}}
\end{frame}

% Slide 11D: Python Laboratory (Terminal Output)
\begin{frame}[fragile]{Python Lab: Terminal Output}
  Exact console execution output verifying numerical precision:
  \vspace{0.2em}
  \begin{block}{Standard Terminal Output}
    \fontsize{6.5pt}{7.5pt}\selectfont
    \begin{verbatim}
--- 1. Frequentist vs Classical Approach (N = 10,000) ---
Coin P(Heads):    Empirical = 0.4987 | Classical = 0.5000
Die  P(X=4):      Empirical = 0.1669 | Classical = 0.1667

--- 2. Axioms & Addition Rule (English Deck) ---
P(Spade):         0.2500 (13/52)
P(King):          0.0769 (4/52)
P(King cap Spade):0.0192 (1/52)
P(Spade cup King):0.3077 (16/52 via Inclusion-Exclusion)
    \end{verbatim}
  \end{block}
\end{frame}

% Slide 16: Closing Reflection
\begin{frame}{Section Closing and Axiomatic Synthesis}
  \begin{reflexion}[The Power of Axioms in Data Science]
    Kolmogorov's three axioms represent a triumph of mathematical synthesis: requiring only non-negativity, normalization, and disjoint additivity, they precisely derive the entire framework of probability theory, safeguarding algorithmic models against logically contradictory inferences.
  \end{reflexion}

  \vspace{0.5em}
  \pause
  \begin{retoclase}[Next Session: Conditional Probability]
    If we learn that a patient tested positive on a diagnostic test ($B$), how do we update the prior probability that they actually have the disease ($A$)? Next session, we introduce the formal mechanism of \emph{Conditional Probability} $P(A \mid B)$.
  \end{retoclase}
\end{frame}

\end{document}
